\BeleskeID{02-s004}
% element: 02-s004:e01
% element: 02-s004:e02
% element: 02-s004:e03
% element: 02-s004:e04
% element: 02-s004:d01

Pretpostavka jednakih kašnjenja služi pojednostavljenom posmatranju pojedinih pojava. Ne znači da će fizičko kolo svaki ulazni talasni oblik samo pomeriti za \(t_p\). Ako su obe izlazne promene ostvarene, invertor na pozitivan ulazni impuls širine \(T_{\mathrm{ul}}=t_2-t_1\) daje negativan impuls širine
\[
T_{\mathrm{izl}}=T_{\mathrm{ul}}+t_{pLH}-t_{pHL}.
\]
Razlika smernih kašnjenja može, dakle, proširiti ili skratiti impuls.

Ovaj račun pretpostavlja da se obe promene zaista prenesu. Kod veoma kratkog impulsa to nije zagarantovano: kolo može početi prelaz, a da novi izlazni nivo ne bude dostignut pre naredne promene ulaza. Izlazni impuls tada može biti izobličen ili izostati. Trajanje ulaznog impulsa zato treba posmatrati u odnosu na odziv kola; samu pretpostavku jednakih kašnjenja ne prenosimo na sve slučajeve.
